% ==============================================================
\section{付録}\label{sec:supplement}
% ==============================================================

\subsection{EX3: パラメータ感度}
\label{sec:supp_ex3}

EX1の語彙内データ（$\beta = 0.3$ 固定）を用いて，
主要パラメータに対する頑健性を評価する（各3 seeds）．

\subsubsection{ウィンドウサイズ $W$}

\begin{table}[h]
\caption{EX3: ウィンドウサイズ感度（3 seeds平均）}
\label{tab:ex3_window}
\centering
\begin{tabular}{ccc}
\toprule
$W$ & F1 & FAR \\
\midrule
10  & 0.75 & 0.67 \\
20  & 0.81 & 0.33 \\
50  & 0.71 & 0.50 \\
100 & 0.62 & 0.33 \\
200 & 0.00 & 0.00 \\
\bottomrule
\end{tabular}
\end{table}

$W = 10$--$100$ で F1 $> 0.6$ を維持する．
$W = 20$ が最高 F1 = 0.81 を達成し，
小さいウィンドウが鋭い変化点検出に有利であることを示す．
$W = 200$ ではサポート変動が平滑化され検出不能となる．

\subsubsection{有意水準 $\alpha$}

\begin{table}[h]
\caption{EX3: 有意水準感度（3 seeds平均）}
\label{tab:ex3_alpha}
\centering
\begin{tabular}{ccc}
\toprule
$\alpha$ & F1 & FAR \\
\midrule
0.01 & 0.00 & 0.00 \\
0.05 & 0.29 & 0.17 \\
0.10 & 0.76 & 0.17 \\
0.20 & 0.71 & 0.50 \\
0.30 & 0.71 & 0.50 \\
\bottomrule
\end{tabular}
\end{table}

$\alpha \geq 0.10$ で安定した検出が得られる．
$\alpha = 0.10$ が最良のF1/FARバランス（F1 = 0.76, FAR = 0.17）を達成する．
$\alpha \leq 0.05$ では BH 補正後の棄却閾値が厳しくなり，
真の帰属も棄却されるため F1 が急落する．

\subsubsection{置換回数 $B$}

\begin{table}[h]
\caption{EX3: 置換回数感度（3 seeds平均）}
\label{tab:ex3_perm}
\centering
\begin{tabular}{ccc}
\toprule
$B$ & F1 & 実行時間 (ms) \\
\midrule
50    & 0.71 & 913 \\
100   & 0.71 & 914 \\
500   & 0.71 & 941 \\
1{,}000 & 0.71 & 969 \\
5{,}000 & 0.71 & 1{,}201 \\
\bottomrule
\end{tabular}
\end{table}

全 $B$ 値で F1 = 0.71 と完全に安定しており，
置換回数に対して頑健である．
$p$ 値の分解能（$1/(B+1)$）を確保するため $B = 5{,}000$ を
デフォルトとする．

\subsubsection{変化点検出法}

閾値交差法とCUSUMを比較した結果，
語彙内ブースト条件では両手法ともF1 = 0.71で同等の精度を示した．

\subsection{EX4: スケーラビリティ}
\label{sec:supp_ex4}

\subsubsection{データサイズ $N$}

\begin{table}[h]
\caption{EX4: データサイズ vs 実行時間}
\label{tab:ex4_n}
\centering
\begin{tabular}{rrrr}
\toprule
$N$ & Phase 1 (s) & Attribution (s) & Total (s) \\
\midrule
1K   & 0.2  & 0.01  & 0.2 \\
5K   & 0.9  & 0.1   & 1.0 \\
10K  & 1.8  & 0.3   & 2.1 \\
50K  & 9.5  & 2.1   & 11.6 \\
100K & 19.2 & 5.1   & 24.4 \\
500K & 102  & 53    & 155 \\
1M   & 217  & 154   & 370 \\
\bottomrule
\end{tabular}
\end{table}

Phase 1 と Attribution の両方がデータサイズに対して
\textbf{線形にスケール}する．
100万トランザクションで約6.2分であり，実用的な処理時間である．

\subsubsection{イベント数 $|\mathcal{E}|$}

\begin{table}[h]
\caption{EX4: イベント数 vs 性能}
\label{tab:ex4_events}
\centering
\begin{tabular}{cccc}
\toprule
$|\mathcal{E}|$ & Attribution (ms) & Total (s) & F1 \\
\midrule
1  & 7    & 0.9 & 1.00 \\
3  & 28   & 0.9 & 1.00 \\
5  & 57   & 1.0 & 0.80 \\
10 & 128  & 1.1 & 0.91 \\
20 & 678  & 1.8 & 0.36 \\
\bottomrule
\end{tabular}
\end{table}

帰属時間はイベント数に対して線形にスケールする．
$|\mathcal{E}| \leq 10$ では F1 $\geq 0.80$ を維持するが，
$|\mathcal{E}| = 20$ では $N = 5{,}000$ に20イベントが密集するため
円形シフト検定の弁別力が低下し F1 = 0.36 となる．
